% Round 15. Formal extension of the separately reviewed P1 integration.
% Requires independent review of the quotient and iteration below.
\section{PWS integration on projective spaces and projective bundles}
\label{r15ps:section}

This section derives a further specific form-valued pushforward
from Theorem~\ref{r15hp:projection} and its algebraic base-change
statement.  The input remains PWS and the prescribed comparison
includes every ordinary smooth compact complementary test.
An explicit tensor projection lemma below also permits compact
PWS tests, still in delta codimension zero.
Neither an all-actual-test proper pushforward nor a pushforward
for arbitrary proper morphisms is asserted.

\begin{lemma}[Projection against a PWS target tensor]
\label{r15ps:PWS-projection-line}
For the projective-line operator of
Theorem~\ref{r15hp:projection}, let $\omega$ be a PWS source
form and let $\beta$ be any PWS target form.  In every meaningful
bidegree one has an equality of ordinary actual PWS forms
\[
 \mathcal I_U(\omega\wedge\pi^*\beta)
                         =(\mathcal I_U\omega)\wedge\beta.
\tag{PS15.0}\label{r15ps:PWS-projection-formula-line}
\]
This assertion is restricted to delta codimension zero in both
inputs.  It is not inferred by inserting a nonsmooth test into
the all-smooth pairing theorem.
\end{lemma}

\begin{proof}
Work near an arbitrary target point.  Shrink so that $\beta$
has one ordinary finite unit-moment PWS presentation.  Perform
the complete-fiber construction for $\omega$ in the proof of
Theorem~\ref{r15hp:projection}.  On its finite splitting cover,
append all pullbacks of the moments of $\beta$ to the base
tuple $b$.  These extra moments are analytic units on the
working ordinary open; they need not be algebraic units on the
whole splitting variety.  This is allowed by
Proposition~\ref{r15ray:analytic}.  On a compact working chart
their enlarged moment image is finite polyhedral.  The source
tensor and the target tensor now have one common finite
presentation in $(b,z-a_i)$ and $b$, respectively.

Use the same root and infinity cutoff $\theta$ for both
$\omega'$ and $\omega'\wedge\pi'^*\beta'$.  It is permitted
for the latter as well: on a discarded tube the first factor
is a base pullback by \eqref{r15hp:graph-difference}, so the
product is also a base pullback, and its wedge with its own
complementary test exceeds the base dimension exactly as in
\eqref{r15hp:tail-degree}.  Its retained moment support is
contained in that of $\theta\omega'$.  Thus both outputs may
be computed by the same finite root-ray bands and the same
finite cutoff bounds.

On any such band write $\Theta$ for the retained source PWS
tensor and $B$ for the target PWS tensor.  Refine by the finitely
many pieces of $B$ as well.  Both vertical tangent vectors
annihilate the pulled-back target tensor.  Therefore the double
vertical contraction has the exact identity
\[
 \operatorname{vc}_t(\Theta\wedge\operatorname{pr}^*B)
                         =\operatorname{vc}_t(\Theta)\wedge B.
\]
There is no additional sign: in the derivation rule the terms
contracting $B$ vanish, and the two remaining contractions
act on the first factor in the same order on both sides.
The coefficients of $B$ depend only on the base variable, so
coefficientwise band integration commutes with this wedge
product.  Tangential restriction to every common face also
commutes with contraction and wedge.  Lemma~\ref{r15ray:band}
therefore gives equality as PWS tensors on the entire base
image, including lower-dimensional faces and collapsing bands.
Ordinary pullback through $b$ proves the desired equality on
the splitting cover.  If a source slot has degree zero, the
corresponding vertical contraction is zero on both sides;
this also covers the zero relative-degree convention.

Let $a$ have native degree $e$.  Both downstairs outputs may
be constructed with this same finite cover.  Apply the finite
PWS trace and divide by $e$.  The frozen finite PWS projection
identity gives
\[
 \operatorname{Tr}^{\mathrm{PS}}_a(\tau'\wedge a^*\beta)
                       =(\operatorname{Tr}^{\mathrm{PS}}_a\tau')
                                                    \wedge\beta.
\]
Its hypotheses are satisfied: both factors are PWS, and in
its pairing proof $\beta\wedge\eta$ for a smooth test $\eta$
is a compact complementary \emph{actual} test, allowed by the
already established finite all-actual-test trace theorem.
This is an existing finite projection identity, not a
nonfinite one being assumed.  The normalized trace construction
and current uniqueness now give
\eqref{r15ps:PWS-projection-formula-line} downstairs.  Ordinary
gluing proves the assertion everywhere.
\end{proof}

\begin{lemma}[The ordered-factor finite quotient]
\label{r15ps:symmetric-quotient}
For every integer $m\geq1$, multiplication of homogeneous binary
linear forms defines a finite surjective morphism
\[
 q_m:(\mathbf P^1_K)^m\longrightarrow\mathbf P^m_K
\tag{PS15.1}\label{r15ps:quotient}
\]
of native generic degree $m!$.  Both source and target are smooth
geometrically integral.  For every smooth integral $K$-variety
$S$, the product map
\[
 Q_{m,S}=\mathrm{id}_S\times q_m:
 S\times(\mathbf P^1)^m\longrightarrow S\times\mathbf P^m
\tag{PS15.2}\label{r15ps:base-quotient}
\]
is finite surjective of the same generic degree $m!$, with
smooth integral source and target.
\end{lemma}

\begin{proof}
Send $([u_1:v_1],\ldots,[u_m:v_m])$ to the coefficients of
$\prod_{j=1}^m(u_jX+v_jY)$.  These multihomogeneous coefficients
never vanish simultaneously.  Thus they define a projective
morphism over the original field $K$.
Over an algebraically closed extension, every nonzero binary
form is a product of linear factors, uniquely as an unordered
multiset with multiplicities.  The geometric fibers of the map
are therefore nonempty and finite.  Properness and
quasi-finiteness make the map finite.  Over the nonempty open
of forms with distinct roots, the ordered factorizations form
a free transitive set under the symmetric group of order $m!$.
Characteristic zero implies separability, so this is the
generic degree.  The asserted smoothness and geometric
integrality are those of products of projective spaces.

Finiteness and surjectivity survive the indicated base change.
At the generic point of $S$, the same distinct-root computation
over $K(S)$ gives degree $m!$.  The products with geometrically
integral smooth projective spaces are smooth integral over
the given integral smooth $S$.  Thus no constant-field factor
is added to or removed from the degree in
\eqref{r15ps:base-quotient}.
\end{proof}

\begin{proposition}[Iteration on a product of projective lines]
\label{r15ps:iterated-lines}
Let $U$ be an ordinary open of a smooth integral algebraic
$K$-variety $S$ of dimension $d$, and let
$p_m:U\times(\mathbf P^1)^m\to U$ be the projection.
Successive application of Theorem~\ref{r15hp:projection}
gives, for every PWS input $\omega$ of bidegree $(p,q)$,
an ordinary PWS form
\[
 \mathcal I^{[m]}_U\omega\in
                    \operatorname{PS}^{p-m,q-m}(U).
\tag{PS15.3}\label{r15ps:iterated-output}
\]
It has the prescribed pairing
\[
 \int_U\mathcal I^{[m]}_U\omega\wedge\eta
      =\int_{U\times(\mathbf P^1)^m}\omega\wedge p_m^*\eta
\tag{PS15.4}\label{r15ps:iterated-pairing}
\]
against every ordinary smooth compact complementary test.
It is independent of the order of the $m$ integrations and
commutes with ordinary restriction and with algebraic base
change between smooth integral varieties.  A negative output
degree means the zero current and the zero form.
\end{proposition}

\begin{proof}
For $p,q\geq m$, integrate one projective-line coordinate at
a time.  At every stage the remaining base is an ordinary open
of a smooth algebraic product of $S$ with projective lines, and
the input produced at the preceding stage is PWS.
Thus every application of Theorem~\ref{r15hp:projection} has
its stated hypotheses.  The final degree is $(p-m,q-m)$.

Pair the last output with a test $\eta$ on $U$.  The pairing
theorem expresses it as an integral on one additional projective
line with the pullback of $\eta$.  That pullback is still an
ordinary smooth test, and is compactly supported because the
projection is proper.  Repeat at each stage.  This proves
\eqref{r15ps:iterated-pairing} for the original arbitrary test,
without expressing it as an exterior product.  Any integration
order gives that same current.  Actual-current injectivity
therefore proves independence of the order.

For an algebraic map $b:T\to S$ of smooth integral varieties,
apply Proposition~\ref{r15hp:base-change} at each stage to the
algebraic map $b$ times the identity on the remaining projective
lines.  Ordinary opens are restricted to the relevant inverse
images.  Induction proves actual base-change compatibility
of the whole iterated operator.  Restriction is proved in the
same way.  If $p<m$ or $q<m$, a negative output degree has no
nonzero complementary test on the $d$-dimensional target;
equivalently the iteration reaches a zero relative-degree
pushforward.  The claimed value is zero.
\end{proof}

\begin{theorem}[PWS integration along a projective space]
\label{r15ps:projective-space}
For $S,U,d$ as in the preceding proposition and every $m\geq0$,
let $\pi_m:U\times\mathbf P^m\to U$ be the projection.
For every $\omega\in\operatorname{PS}^{p,q}
(U\times\mathbf P^m)$ its prescribed proper current has the
unique ordinary PWS representative
\[
 \mathcal I_{\pi_m}\omega
   =\frac1{m!}\mathcal I^{[m]}_U(Q_{m,S}^*\omega)
       \in\operatorname{PS}^{p-m,q-m}(U),
\tag{PS15.5}\label{r15ps:space-output}
\]
with the evident identity convention for $m=0$.
For every ordinary smooth compact complementary test $\eta$,
\[
 \int_U\mathcal I_{\pi_m}\omega\wedge\eta
       =\int_{U\times\mathbf P^m}\omega\wedge\pi_m^*\eta.
\tag{PS15.6}\label{r15ps:space-pairing}
\]
This operator commutes with ordinary restriction and with every
algebraic base change between smooth integral varieties.
\end{theorem}

\begin{proof}
The finite map $Q_{m,S}$ is a genuine algebraic map, so its
analytic pullback preserves PWS forms.  Proposition~\ref{r15ps:iterated-lines}
makes the right side of \eqref{r15ps:space-output} an ordinary
PWS form of the asserted degree.
Its pairing with an arbitrary test is
\[
 \begin{split}
 \int_U\mathcal I_{\pi_m}\omega\wedge\eta
 &=\frac1{m!}\int_{U\times(\mathbf P^1)^m}
                         Q_{m,S}^*\omega\wedge p_m^*\eta\\
 &=\frac1{m!}\int_{U\times(\mathbf P^1)^m}
                      Q_{m,S}^*(\omega\wedge\pi_m^*\eta)\\
 &=\int_{U\times\mathbf P^m}\omega\wedge\pi_m^*\eta.
 \end{split}
\tag{PS15.7}\label{r15ps:one-native-factor}
\]
The final step is exactly the native finite degree formula
for $Q_{m,S}$ of degree $m!$.  The paired input is an actual
PWS top form of compact support.  There is a single finite
quotient and a single factor $m!$ in this calculation; neither
the iteration nor a change of presentation introduces another
degree factor.  The formula remains the native one if the
integral variety $S$ has a larger constant field than $K$:
source and target have that same constant field, and their
relative generic degree is still $m!$.

Current injectivity gives uniqueness.  Restriction follows from
the formula.  For algebraic base change $b:T\to S$, the finite
maps $Q_{m,T}$ and $Q_{m,S}$ form a Cartesian product square.
Functoriality of genuine analytic pullback gives
\[
 Q_{m,T}^*(b\times\mathrm{id}_{\mathbf P^m})^*\omega
       =(b\times\mathrm{id}_{(\mathbf P^1)^m})^*
                                                     Q_{m,S}^*\omega.
\]
Apply the iterated base-change identity of
Proposition~\ref{r15ps:iterated-lines} and divide by the same
$m!$.  This proves actual naturality.  It uses no base-change
identity for a finite trace or a nonfinite current.
\end{proof}

\begin{theorem}[PWS integration along a projective bundle]
\label{r15ps:projective-bundle}
Let $S/K$ be smooth and integral, let $N$ be an algebraic vector
bundle of rank $m+1$, and put
$\pi:E=\mathbf P_S(N)\to S$.  On any ordinary open
$U\subset S^{\mathrm{an}}$, every
$\omega\in\operatorname{PS}^{p,q}(E_U)$ has a unique ordinary
PWS pushforward
\[
 \mathcal I_\pi\omega\in\operatorname{PS}^{p-m,q-m}(U)
\tag{PS15.8}\label{r15ps:bundle-output}
\]
with prescribed pairing
\[
 \int_U\mathcal I_\pi\omega\wedge\eta
                    =\int_{E_U}\omega\wedge\pi^*\eta
\tag{PS15.9}\label{r15ps:bundle-pairing}
\]
for every ordinary smooth compact complementary test.
It commutes with ordinary restriction and with pullback of
the vector bundle by every algebraic map between smooth
integral bases.  It satisfies the projection formula against
ordinary PWS target forms.  If the input is killed by
$d'$, $d''$, or $D$, then its output is killed by that same
operator as an actual form.
\end{theorem}

\begin{proof}
Choose Zariski opens of $S$ on which $N$ is trivial.  They give
ordinary opens $U_i$ covering $U$ and algebraic identifications
$E_{U_i}\simeq U_i\times\mathbf P^m$.
Theorem~\ref{r15ps:projective-space} produces an ordinary PWS
form on each $U_i$.  Under any two trivializations it has the
same prescribed current \eqref{r15ps:bundle-pairing} on every
ordinary overlap: the identifications are actual isomorphisms,
whose change-of-variables degree is one.  Current injectivity
makes the forms equal.  They glue in the ordinary topology.
Smooth partitions of unity on the target extend the local
pairing equality to every compact test on $U$.  Properness of
$\pi$ keeps all pulled-back paired supports compact.

If $b:T\to S$ is algebraic between smooth integral bases, its
pullback bundle is $E_T=\mathbf P_T(b^*N)$, and the pulled
trivializations identify its Cartesian diagram locally with
the product base-change diagram in
Theorem~\ref{r15ps:projective-space}.  That theorem proves
\[
 b^*\mathcal I_\pi\omega
        =\mathcal I_{\pi_T}(b_E^*\omega)
\tag{PS15.10}\label{r15ps:bundle-base-change}
\]
on ordinary opens where the expressions are defined.  Ordinary
gluing gives the global assertion.

For a PWS target form $\nu$, first apply
Lemma~\ref{r15ps:PWS-projection-line} at every stage of the
iterated projective-line integration.  At each stage the
remaining target factor is the pullback of $\nu$, and is PWS.
The iterated operator therefore commutes with its wedge.
The finite map $Q_{m,S}$ pulls back
$\omega\wedge\pi_m^*\nu$ to
$Q_{m,S}^*\omega\wedge p_m^*\nu$.  Formula
\eqref{r15ps:space-output} proves the same projection identity
for a projective space.  Applying this on every trivializing
ordinary open and gluing gives
\[
 \mathcal I_\pi(\omega\wedge\pi^*\nu)
                           =(\mathcal I_\pi\omega)\wedge\nu
\tag{PS15.11}\label{r15ps:PWS-projection}
\]
as an actual PWS identity.  This deduction uses the explicit
tensor lemma, not an insertion of a nonsmooth test into
\eqref{r15ps:bundle-pairing}.

Finally proper pushforward of currents commutes with the full
actual differentials, as follows directly by actual Stokes
with differentiated compact smooth tests.  The relative real
degree is $2m$, so no odd-degree sign is introduced.  Thus, if
the corresponding full derivative of $\omega$ is zero, that
of the actual output has zero current and is zero by current
injectivity.  This last assertion does not apply the PWS
pushforward theorem to a derivative with positive delta
codimension.  It uses only the current identity and the fact
that the derivative of the already constructed output is an
actual form.
\end{proof}

\begin{corollary}[Compact PWS test pairings]
\label{r15ps:PWS-tests}
For the projective-line, projective-space, and projective-bundle
operators above, the prescribed pairing also holds against
every compactly supported complementary ordinary PWS test.
No test with positive delta codimension is included by this
corollary.
\end{corollary}

\begin{proof}
Let $\eta$ be such a PWS test on the base.  The source form
$\omega\wedge\pi^*\eta$ is PWS and compactly supported,
since $\pi$ is proper.  The projection identity gives the
top-degree ordinary PWS output
\[
 \mathcal I_\pi(\omega\wedge\pi^*\eta)
                             =(\mathcal I_\pi\omega)\wedge\eta.
\]
Its support lies in $\operatorname{supp}\eta$ by ordinary
restriction compatibility.  Choose an ordinary smooth compact
scalar $\zeta$ equal to one on a neighborhood of
$\operatorname{supp}\eta$.  The original all-smooth pairing
theorem, now applied to the PWS source
$\omega\wedge\pi^*\eta$ and the smooth scalar test $\zeta$,
gives
\[
 \int_U(\mathcal I_\pi\omega)\wedge\eta
     =\int_{E_U}\omega\wedge\pi^*\eta.
\]
The same choice of $\zeta$ makes it drop out of the source
integral.  This proves the assertion without extending the
input or test category beyond PWS.
\end{proof}

\paragraph{Boundary of the new interface.}
The theorem constructs PWS outputs for projective-bundle
projections of any relative dimension, with all ordinary
smooth-test pairings and the stated algebraic base change.
It does not itself compute Aeppli cohomology of a nontrivial
projective bundle.  In particular, a higher-rank relative
diagonal is not a divisor, so the rank-two scalar diagonal
kernel does not automatically extend.  Any incidence or flag
argument using this interface must separately justify its
cohomological identities, choice of smooth Chern forms,
primitive degrees, and exceptional quotient.
